\label{sec:pas}

\textbf{This section is illustrative, not the main claim.}
Section~\ref{sec:taxonomy} argues for a missing safety layer; here we show only
that the principles are \emph{implementable} on a physical platform
(SocialVLA-Engage, running at 10\,Hz) using \textbf{PAS}
(\emph{Probe--Authorize--Speak}).

\paragraph{Runtime data flow.}
Each control cycle: perception updates engagement cues
\citep{rich2010,salam2017}; the \emph{probe} module emits only reversible
nonverbal actions (\texttt{wait}, \texttt{head}, \texttt{orient}) and never
\texttt{speak}; evidence accumulates from probe responses
\citep{skantze2014speech,bohus2014}; the \emph{authorization gate} is the sole
speak path; only after authorization does a foundation-model dialogue stack run.
Passive direct-init ($e_t > \theta$) bypasses this split and can speak when
refraining is socially safer.
Any architecture enforcing probe $\rightarrow$ authorize $\rightarrow$ act would
serve the same pedagogical role.

\begin{algorithm}[!htb]
\caption{Initiation authorization loop (10\,Hz sketch)}
\label{alg:pas}
\small
\begin{algorithmic}[1]
\STATE \textbf{Input:} cues $e_t$, probe history, venue risk posture
\STATE \textbf{Invariant:} probe module $\not\rightarrow$ \texttt{speak}
\IF{authorization gate closed}
  \STATE emit reversible probe or \texttt{wait}; update receptivity evidence
\ELSIF{gate open}
  \STATE permit FM dialogue / first \texttt{speak}; log decision margin
\ENDIF
\end{algorithmic}
\end{algorithm}

\subsection{Evaluation Protocol}

This is a \textbf{position and design submission}.
Quantitative outcomes from an ongoing between-subjects study ($n{=}90$, three
policy classes, ethics-approved) are \textbf{in progress}.
Rethinking Safety explicitly welcomes in-progress work; our goal is not
placeholder effect sizes but whether initiation authorization deserves a place in
generalist-robot safety discourse.
Empirical numbers will test the sketch; they do not define the framing.
All arms share hardware, a 10\,Hz perception loop, and the same post-initiation
dialogue stack; only initiation policy differs:
\begin{center}
\scriptsize
\begin{tabular}{@{}l p{0.62\linewidth}@{}}
\toprule
\textbf{Condition} & \textbf{Policy class} \\
\midrule
PAS & Probes + authorization gate (default risk posture) \\
Direct-init & $e_t > \theta$; no probes \\
Passive-wait & Probes + conservative gate (rarely speak) \\
\bottomrule
\end{tabular}
\end{center}

\paragraph{Endpoints.}
\textbf{Mechanistic:} distribution of speak-margin proxy at first
\texttt{speak} ($\widehat{\Delta}_{\text{init}}$), comparing PAS vs.\
direct-init.
\textbf{User-reported:} conversation success, awkwardness, and naturalness
(7-point Likert).
\textbf{Deployment:} risk-posture sweep on PAS without retraining perception.

\paragraph{Falsifiers.}
Even before numbers arrive, the framing is testable:
if PAS and direct-init produce indistinguishable margin distributions, the
architectural split adds little;
if awkwardness does not covary with margin alignment, felt social safety may
require metrics beyond timing proxies \citep{tian2021socialerrors};
if risk-posture settings do not monotonically trade initiation rate for margin,
deployment-time authorization is not a viable safety dial.
A single-site study cannot certify generalist-robot safety; it can only show
that initiation authorization is \emph{measurable and manipulable}---enough, we
hope, to motivate the community benchmarks in Section~\ref{sec:open}.
